% ============================================================
% DOCUMENTCLASS
% ============================================================
\documentclass[12pt,oneside]{book}

% ============================================================
% METADATOS
% ============================================================
\newcommand{\universidad}{Universidad de Buenos Aires (UBA)}
\newcommand{\facultad}{}
\newcommand{\carrera}{}
\newcommand{\materia}{Derecho Procesal Civil}
\newcommand{\titulo}{Unidad 7: Actos Procesales}
\newcommand{\autor}{Isaac Edgar Camacho Ocampo}
\newcommand{\anio}{2024}

% ============================================================
% PAQUETES
% ============================================================
\usepackage[utf8]{inputenc}
\usepackage[T1]{fontenc}
\usepackage[spanish,es-nodecimaldot]{babel}

\usepackage[
    top=2.5cm,
    bottom=2.5cm,
    left=3cm,
    right=2.5cm,
    headheight=15pt
]{geometry}

\usepackage{amsmath}
\usepackage{amssymb}
\usepackage{mathtools}

\usepackage{graphicx}
\usepackage{float}
\usepackage{caption}
\usepackage{subcaption}

\usepackage{booktabs}
\usepackage{array}
\usepackage{longtable}
\usepackage{tabularx}

\usepackage{enumitem}

\usepackage{xcolor}
\usepackage[most]{tcolorbox}

\usepackage{fancyhdr}

\usepackage{ragged2e}

\usepackage[
    colorlinks=true,
    linkcolor=blue!50!black,
    citecolor=green!40!black,
    urlcolor=blue!60!black,
    pdftitle={\titulo},
    pdfauthor={\autor},
    pdfsubject={Apuntes universitarios},
    bookmarks=true
]{hyperref}

% ============================================================
% RUTAS DE IMÁGENES
% ============================================================
\graphicspath{
    {./img/}
    {./graficos/}
    {./figuras/}
}

% ============================================================
% CONFIGURACIÓN GENERAL
% ============================================================
\setcounter{tocdepth}{2}

\setlength{\parindent}{0pt}
\setlength{\parskip}{0.5em}

\setlist[itemize]{topsep=0.3em, itemsep=0.2em}
\setlist[enumerate]{topsep=0.3em, itemsep=0.2em}

\clubpenalty=10000
\widowpenalty=10000

% ============================================================
% ENCABEZADOS Y PIES
% ============================================================
\pagestyle{fancy}
\fancyhf{}
\fancyhead[L]{\nouppercase{\leftmark}}
\fancyhead[R]{\materia}
\fancyfoot[C]{\thepage}
\renewcommand{\headrulewidth}{0.4pt}
\renewcommand{\footrulewidth}{0pt}

\fancypagestyle{plain}{
    \fancyhf{}
    \fancyfoot[C]{\thepage}
    \renewcommand{\headrulewidth}{0pt}
}

% ============================================================
% COMANDOS MATEMÁTICOS
% ============================================================
\newcommand{\abs}[1]{\left|#1\right|}
\newcommand{\norm}[1]{\left\lVert#1\right\rVert}
\newcommand{\vect}[1]{\mathbf{#1}}
\newcommand{\deriv}[2]{\frac{d#1}{d#2}}
\newcommand{\pderiv}[2]{\frac{\partial#1}{\partial#2}}

% ============================================================
% ENTORNOS / CAJAS ACADÉMICAS
% ============================================================
\newtcolorbox{definition}[1]{
    enhanced,
    breakable,
    title={Definición: #1},
    fonttitle=\bfseries,
    colback=blue!3,
    colframe=blue!50!black,
    boxrule=0.8pt,
    arc=2mm
}

\newtcolorbox{important}{
    enhanced,
    breakable,
    title={Importante},
    fonttitle=\bfseries,
    colback=yellow!8,
    colframe=orange!70!black,
    boxrule=0.8pt,
    arc=2mm
}

\newtcolorbox{example}[1]{
    enhanced,
    breakable,
    title={Ejemplo: #1},
    fonttitle=\bfseries,
    colback=green!4,
    colframe=green!40!black,
    boxrule=0.8pt,
    arc=2mm
}

\newtcolorbox{formula}[1]{
    enhanced,
    breakable,
    title={Fórmula / Regla: #1},
    fonttitle=\bfseries,
    colback=purple!4,
    colframe=purple!50!black,
    boxrule=0.8pt,
    arc=2mm
}

\newtcolorbox{exercise}[1]{
    enhanced,
    breakable,
    title={Ejercicio: #1},
    fonttitle=\bfseries,
    colback=gray!5,
    colframe=gray!60!black,
    boxrule=0.8pt,
    arc=2mm
}

\newtcolorbox{note}{
    enhanced,
    breakable,
    title={Nota},
    fonttitle=\bfseries,
    colback=white,
    colframe=gray!50,
    boxrule=0.6pt,
    arc=2mm
}

% ============================================================
% INICIO DEL DOCUMENTO
% ============================================================
\begin{document}

% ============================================================
% PORTADA
% ============================================================
\begin{titlepage}

\thispagestyle{empty}

\begin{center}

\vspace*{1cm}

{\large \textsc{\universidad}}

\vspace{3cm}

{\Huge \textbf{\materia}}

\vspace{1cm}

{\LARGE \titulo}

\vspace{0.8cm}

\textit{Comprender los conceptos es el primer paso para convertir el estudio en conocimiento.}

\vspace{1.2cm}

\rule{0.70\textwidth}{0.5pt}

\vspace{0.5cm}

{\large Apuntes de estudio}

\vspace{0.5cm}

\rule{0.70\textwidth}{0.5pt}

\vfill

{\large \autor}

\vspace{0.3cm}

{\large \anio}

\end{center}

\vfill

\begin{flushleft}
\textit{Este es un modesto aporte a los estudiantes de ``Derecho Procesal Civil'' no pretende sustituir el material oficial de la materia.}
\end{flushleft}

\end{titlepage}

% ============================================================
% ÍNDICE
% ============================================================
\tableofcontents
\clearpage

% ============================================================
% CAPÍTULO 1: CONCEPTO Y CLASIFICACIÓN DE LOS ACTOS PROCESALES
% ============================================================
\chapter{Concepto y Clasificación de los Actos Procesales}

\section{Concepto de Acto Procesal}

\begin{definition}{Acto Procesal}
Son actos procesales los \textbf{hechos voluntarios} que tienen por efecto inmediato lograr la \textbf{constitución}, el \textbf{desenvolvimiento} o la \textbf{extinción/conclusión} de un proceso judicial.
\end{definition}

El proceso es una concatenación, una cadena sucesiva de actos procesales que van desde un inicio hasta la conclusión normal, que es la \textbf{sentencia}. La sentencia es el modo típico de terminación, aunque el proceso puede concluir antes de llegar a ella a través de los denominados \textbf{modos anormales de extinción}.

Los modos anormales son: el \textbf{desistimiento}, la \textbf{conciliación}, la \textbf{transacción}, el \textbf{allanamiento} y la \textbf{caducidad de instancia}. Estos se analizarán en detalle en la sección correspondiente.

\section{Distinción: Acto Procesal vs. Hecho Procesal}

En el derecho civil existe la distinción entre hechos jurídicos y actos jurídicos. Análogamente, en el proceso se distingue entre hechos procesales y actos procesales.

\begin{definition}{Hecho Procesal}
Son circunstancias que impactan directamente en el curso del proceso, pero que \textbf{no responden a la voluntad de las partes}, sino a acontecimientos naturales o ajenos a la voluntad humana.
\end{definition}

\begin{example}{Hechos procesales}
La muerte de una de las partes durante el proceso; la restricción de la capacidad o la caída en algún tipo de enfermedad mental, que hace que alguno de los sujetos procesales pierda la plena capacidad y deba ser representado por otro.
\end{example}

Los \textbf{actos procesales}, en cambio, son aquellos que responden a hechos voluntarios, actuados por la persona humana: las partes, sus letrados, el juez, el secretario y todos los demás sujetos procesales.

\section{Sujetos que Realizan Actos Procesales}

\subsection{Las Partes o Peticionarios}

Son los sujetos que están por un lado (la \textbf{parte actora}) o por el otro (la \textbf{parte demandada}) en el litigio. El actor puede ser uno solo; el demandado también. Pero puede haber pluralidad de actores y de demandados en un mismo proceso, en los supuestos de \textbf{litis consorcio}.

Tanto la parte actora como la demandada, y sus distintos litis consortes, son sujetos que pueden realizar actos procesales: actos de petición, actos de impulso, de inicio, de desenvolvimiento o de extinción del proceso.

\subsection{Los Auxiliares de las Partes}

Son los sujetos que trabajan con las partes y las auxilian: fundamentalmente los \textbf{letrados patrocinantes} y los \textbf{procuradores apoderados}. Estos profesionales pueden realizar actos procesales en nombre y representación de las partes, o autónomamente, pero siempre al servicio de la parte que patrocinan.

\subsection{El Órgano Judicial y sus Auxiliares}

El \textbf{juez} realiza actos procesales a través de sus resoluciones, providencias y decisiones, en respuesta a las peticiones de las partes. Sus auxiliares incluyen:

\begin{itemize}
    \item El secretario del juzgado.
    \item El prosecretario administrativo.
    \item Los oficiales y agentes del tribunal.
\end{itemize}

\subsection{Terceros Vinculados al Proceso}

También pueden ser sujetos activos de actos procesales:

\begin{itemize}
    \item Los peritos designados por el tribunal, para asesorar en temas de especialidad técnica.
    \item Los oficiales de justicia.
    \item Los oficiales notificadores.
    \item Los martilleros.
    \item Los testigos.
\end{itemize}

Todos estos son sujetos que colaboran al desarrollo del proceso a través de los actos procesales que tienen la carga de realizar.

\section{Clasificación de los Actos Procesales según Palacio}

La clasificación que se sigue en este material es la propuesta por el Dr. \textbf{Lino Palacio}, considerada de interés particular por su claridad y sistematicidad. Pueden encontrarse otras clasificaciones en diferentes autores.

Los actos procesales se agrupan en tres grandes categorías: \textbf{actos de iniciación}, \textbf{actos de desarrollo} y \textbf{actos de conclusión}.

\section{Actos de Iniciación}

Son los que dan comienzo al proceso. Pertenecen fundamentalmente a la primera etapa del proceso: la \textbf{etapa constitutiva}.

Esta etapa suele comenzar, ordinariamente, con:

\begin{itemize}
    \item Las \textbf{diligencias preliminares}: actos procesales que se realizan antes del inicio formal del proceso, orientados a obtener la información necesaria para poder armar adecuadamente una demanda.
    \item La \textbf{demanda}: donde la parte actora plantea su pretensión.
    \item La \textbf{contestación de demanda}: donde la parte demandada plantea su postura.
\end{itemize}

\section{Actos de Desarrollo}

Son los actos de desenvolvimiento del proceso, una vez que ya le fue dado inicio. Se dividen en dos grandes categorías: \textbf{actos de instrucción} y \textbf{actos de dirección}.

\subsection{Actos de Instrucción}

Se asocian con los hechos. En el proceso civil, las partes son las que aportan al juez los hechos: a través del relato que formulan y de las pruebas que permiten reconstruirlos.

\begin{definition}{Actos de Alegación}
Son aquellos por los cuales las partes incorporan los hechos al proceso a través de su relato. Puede hacerse en la demanda, al plantearse un hecho nuevo, o al incorporarse un hecho sobreviniente.
\end{definition}

\begin{definition}{Actos de Prueba}
Son aquellos orientados a reconstruir los hechos alegados, aportando los elementos que permiten al juez verificarlos.
\end{definition}

\subsection{Actos de Dirección}

Se dividen en cuatro categorías:

\begin{table}[H]
\centering
\begin{tabularx}{\textwidth}{>{\bfseries\raggedright\arraybackslash}p{0.22\textwidth} >{\raggedright\arraybackslash}X}
\toprule
\textbf{Tipo} & \textbf{Descripción} \\
\midrule
Actos de Ordenación & Tienden a organizar el proceso y darle un orden. Se subdividen en: actos de impulso, actos de resolución y actos de impugnación. \\[0.4em]
Actos de Transmisión & Disponen poner en conocimiento de las partes las decisiones adoptadas o los pedidos formulados. Ejemplo: traslados. \\[0.4em]
Actos de Documentación & Incorporan al proceso elementos documentales de prueba en cualquier soporte (papel, filmación, informes de organismos, etc.). \\[0.4em]
Actos Cautelares & Orientados a conservar situaciones de hecho o patrimoniales para asegurar que la sentencia pueda ejecutarse. \\
\bottomrule
\end{tabularx}
\caption{Subcategorías de los actos de dirección}
\end{table}

\subsubsection{Actos de Ordenación: Actos de Impulso}

Son los actos que hacen \textbf{avanzar el proceso} hacia su destino ordinario, que es la sentencia. Son todos aquellos actos que van haciendo pasar la cadena de un eslabón al otro, permitiendo que el proceso progrese.

\begin{important}
Los actos de impulso son los únicos que \textbf{interrumpen el plazo de caducidad de instancia}. Esta categoría tiene por ello una relevancia práctica fundamental.
\end{important}

\begin{example}{Qué NO es un acto de impulso}
Que uno de los litigantes cambie de abogado y se presente un nuevo letrado en alguna etapa del proceso: es un acto relevante (cambia el domicilio, puede cambiar la estrategia), pero no hace avanzar el proceso. Tampoco lo es que una parte justifique su inasistencia a una audiencia por tener un viaje: es relevante que el juez lo sepa, pero no impulsa nada.
\end{example}

\subsubsection{Actos de Ordenación: Actos de Resolución o Decisión}

Son los actos que dispone el \textbf{tribunal} en respuesta a las peticiones que hacen las partes. Las partes piden alguna actuación del tribunal; el tribunal resuelve, decide si se hace lugar o no, expresando las razones correspondientes.

\subsubsection{Actos de Ordenación: Actos de Impugnación}

Son los actos que buscan \textbf{modificar o reemplazar} una actuación judicial o actos procesales anteriores, por la vía de los recursos o del incidente de nulidad.

Los recursos fundamentales son:
\begin{itemize}
    \item El \textbf{recurso de revocatoria o reposición}: busca modificar una resolución judicial a través del mismo órgano que la dictó.
    \item El \textbf{recurso de apelación}: busca modificar una resolución a través del órgano superior, el tribunal de segunda instancia.
    \item La \textbf{aclaratoria}: no se categoriza exactamente como recurso, pero cumple una función análoga en términos de revisión.
\end{itemize}

El \textbf{incidente de nulidad} busca que se decrete la nulidad de un acto por algún vicio o defecto en su realización ---un acto cumplido en violación a alguna pauta normativa---, anulando sus efectos y los de los actos ulteriores que se tomaron como consecuencia.

\subsubsection{Actos Cautelares}

Las \textbf{medidas cautelares} están orientadas a conservar situaciones de hecho o patrimoniales para evitar que, al momento del dictado de la sentencia, no haya material para tornar realizable la condena.

\begin{example}{Medida de no innovar}
Una medida de no innovar busca mantener una situación de hecho y evitar que se modifique. Si el bien que se pretende recuperar desaparece del patrimonio del demandado (porque lo enajenó, donó o se deshizo de él), la sentencia no podrá cumplirse. Las medidas cautelares previenen esto.
\end{example}

\section{Actos de Conclusión del Proceso}

El acto prototípico de conclusión del proceso es la \textbf{sentencia}, que es el modo normal de terminación. Pero existen otros modos denominados anormales.

\subsection{Modos Anormales de Terminación del Proceso}

\begin{table}[H]
\centering
\begin{tabularx}{\textwidth}{>{\bfseries\raggedright\arraybackslash}p{0.22\textwidth} >{\raggedright\arraybackslash}X >{\raggedright\arraybackslash}p{0.22\textwidth}}
\toprule
\textbf{Modo} & \textbf{Descripción} & \textbf{Voluntad} \\
\midrule
Desistimiento & La parte actora decide no seguir adelante con el juicio y abandona su pretensión voluntariamente. & Unilateral (actor) \\[0.4em]
Allanamiento & La parte demandada se somete unilateralmente a la pretensión de la actora. & Unilateral (demandado) \\[0.4em]
Conciliación & Ambas partes logran un avenimiento/acuerdo dentro del proceso, generalmente en una audiencia con concesiones recíprocas. & Bilateral (en el proceso) \\[0.4em]
Transacción & Contrato extrajudicial por el cual las partes se hacen concesiones recíprocas para poner fin a los derechos litigiosos. Se plasma fuera del proceso y se presenta al juicio. & Bilateral (fuera del proceso) \\[0.4em]
Caducidad de Instancia & Transcurrido el plazo legal sin actos de impulso, el ordenamiento presume abandono del proceso. Es al mismo tiempo una sanción procesal. & Presunta (inactividad) \\
\bottomrule
\end{tabularx}
\caption{Modos anormales de terminación del proceso}
\end{table}

\begin{formula}{Art. 310 CPCCN --- Caducidad de Instancia}
Caducidad de instancia: en el \textbf{juicio ordinario} el plazo es de \textbf{seis meses}; en los \textbf{procesos de ejecución e incidentes}, de \textbf{tres meses}. Transcurrido ese tiempo sin que las partes hayan realizado actos de impulso, el ordenamiento presume abandono del proceso y habilita su extinción.
\end{formula}

% ============================================================
% CAPÍTULO 2: FORMA, LUGAR Y TIEMPO DE LOS ACTOS PROCESALES
% ============================================================
\chapter{Forma, Lugar y Tiempo de los Actos Procesales}

\section{La Forma de los Actos Procesales}

La \textbf{forma} es la manera en que el acto procesal se exterioriza y se vuelca en el proceso; la manera en que el acto procesal debe ser cumplido. Hay dos dimensiones: el \textbf{modo de expresión} y el \textbf{modo de recepción e incorporación al proceso}.

\subsection{Modo de Expresión: el Lenguaje}

El ordenamiento establece que el idioma debe ser el \textbf{idioma nacional}. Los actos realizados en el extranjero o en otro idioma pueden incorporarse al proceso con su debida traducción realizada por traductor público nacional.

En el caso de personas que no pueden hablar en idioma nacional por ser sordomudas, se designa un \textbf{intérprete del idioma de señas}, debidamente registrado y habilitado, para que interprete y pase al idioma nacional escrito lo que la persona expresa con el lenguaje de sus manos.

\subsection{Características de los Escritos Judiciales}

Hay una regulación detallada a través de acordadas de la Corte que establece cómo deben presentarse los escritos judiciales:

\begin{itemize}
    \item Deben estar escritos en \textbf{tinta negra} (no roja, no verde, sin subrayados en color).
    \item Deben contener una \textbf{SUMA}: título que resume sintéticamente el contenido del escrito.
\end{itemize}

\begin{example}{Ejemplos de sumas}
`Impugna pericia médica' --- `Apela sentencia' --- `Presenta memorial' --- `Opone excepciones'. Con solo leer la suma se sabe de qué trata el escrito.
\end{example}

Todo escrito judicial debe contener en su \textbf{encabezamiento}:

\begin{itemize}
    \item Nombre de quien formula la petición.
    \item Nombre del letrado patrocinante o apoderado (con tomo y folio ante el Colegio Público de Abogados).
    \item \textbf{Carátula del expediente}: nombre del actor, nombre del demandado y objeto del juicio.
    \item \textbf{Domicilio electrónico} (que hoy reemplaza al domicilio procesal físico), donde son válidas todas las notificaciones fehacientes del proceso.
\end{itemize}

\begin{example}{Ejemplo de carátula}
`Carlos O. Funes contra Empresa de Transporte La Nueva Ideal S.A. sobre Daños y Perjuicios'
\end{example}

\subsection{Notas Manuscritas (Art. 117 CPCCN)}

\begin{formula}{Art. 117 CPCCN}
Se admite la nota manuscrita en el expediente para peticiones de mero trámite. Esta nota se trata igual que un escrito: se le coloca el cargo como si fuera una presentación formal.
\end{formula}

\subsection{Presentación de Copias --- Art. 120 CPCCN (Régimen Actual)}

\begin{formula}{Art. 120 CPCCN --- Copias}
De todo escrito judicial del que deba darse traslado, la parte que lo presenta debe acompañar tantas copias como partes contrarias haya. Hoy, esta exigencia se cumple de manera virtual: dentro de las \textbf{24 horas} de presentado el escrito en la mesa de entradas, el abogado debe subir el scan o archivo al sistema de gestión judicial (\textbf{NextGen}).
\end{formula}

\begin{important}
La omisión en el cumplimiento de ese recaudo puede hacer \textbf{tener por no presentado el escrito}, lo que puede significar un gravísimo perjuicio para el cliente. El expediente se manda a desglozar y se lo tiene por no presentado.
\end{important}

\section{Modo de Recepción: El Cargo}

Los escritos se reciben en el tribunal con la colocación del \textbf{cargo}, que es una constancia de recepción que establece:

\begin{itemize}
    \item El dato del tribunal.
    \item La fecha en que se hizo la presentación.
    \item Si la presentación vino firmada por letrado o sin firma de letrado.
\end{itemize}

El cargo es suscripto por el \textbf{prosecretario administrativo} (función típica), aunque puede ser firmado por el secretario o incluso el juez. Es lo que le da la \textbf{fecha cierta} al escrito judicial y lo transforma en un acto público incorporado al expediente judicial.

\section{Las Audiencias: Actas y Videofilmación}

Las audiencias pueden incorporarse al proceso mediante:

\begin{itemize}
    \item \textbf{Actas escritas}: folios redactados por el secretario o el juez donde se vuelca un resumen de lo acontecido, con las firmas de los comparecientes.
    \item \textbf{Videofilmación}: incorporación al sistema de gestión judicial (NextGen) de un archivo de video.
\end{itemize}

Hoy se usan ambos complementariamente: el acta documenta la comparecencia y firmas; el video contiene el desarrollo completo. El sistema de gestión permite acceder a las audiencias videofilmadas de manera remota desde el estudio jurídico, sustituyendo en la práctica la antigua acta judicial escrita.

\subsection{Publicidad de las Audiencias}

En principio, las audiencias son \textbf{públicas}: cualquier persona puede concurrir y presenciar lo que ocurre (sin participar, salvo que el juez lo autorice).

\textbf{Excepciones:} asuntos de familia (siempre reservados por involucrar la intimidad familiar e intereses de menores) y casos donde las partes lo pidan o el juez lo disponga por razones fundadas.

\subsection{Presencia del Juez en Audiencias}

\begin{itemize}
    \item \textbf{Audiencia preliminar (art. 360):} debe ser presidida por el juez \textbf{bajo pena de nulidad}. En ella también suele tomarse la audiencia confesional.
    \item \textbf{Audiencias de prueba} (testimoniales, etc.): pueden ser presenciadas por el juez o \textbf{delegadas} al secretario del tribunal.
\end{itemize}

\begin{formula}{Art. 125, inc. 3°, CPCCN}
Las audiencias deben ser fijadas con una antelación no menor a los \textbf{tres días}, salvo que por razones de urgencia el juez establezca excepcionalmente un plazo menor. La parte que no concurra pierde los derechos que sólo puede ejercer con su presencia.
\end{formula}

\section{Lugar de los Actos Procesales}

El lugar prototípico donde se llevan a cabo los actos procesales es la \textbf{sede del tribunal} ---el juzgado---. La mayoría de los actos ocurre allí.

\subsection{Excepciones: Actos Fuera de la Sede del Tribunal}

\begin{enumerate}[label=\alph*)]
    \item \textbf{Audiencias en el domicilio del declarante:} cuando una parte o testigo no puede trasladarse al tribunal por razones de salud, edad u otras circunstancias justificadas, el tribunal puede constituirse en el lugar donde se encuentra la persona.

    \item \textbf{Reconocimiento judicial de lugares:} el tribunal debe trasladarse físicamente al lugar cuando ello resulta necesario.

    \begin{example}{Street View como herramienta judicial}
    Un accidente de tránsito ocurrido en la Avenida 9 de Julio y Avenida Santa Fe en 2011: hoy concurrir al lugar sería inútil porque la dinámica del tráfico cambió radicalmente (en 2011 no existía el Metrobús). En cambio, con Street View es posible acceder a las imágenes del lugar tal como estaba en 2011, lo cual brinda una reconstrucción más fiel que la inspección ocular actual. Esta herramienta tiene hoy amplia aceptación en los tribunales.
    \end{example}

    \item \textbf{Actos en extraña jurisdicción mediante Oficio Ley 22.172:} cuando debe realizarse una diligencia en una provincia diferente.

    \begin{example}{Oficio Ley 22.172}
    Un juez de Capital Federal necesita constatar el estado de plantaciones de arándanos en un campo de Entre Ríos y ordenar su allanamiento. No tiene jurisdicción para disponer esa medida allí, de modo que libra un Oficio Ley 22.172 al juez de Entre Ríos, quien ordena a la policía y al oficial de justicia del lugar constituirse en el campo, realizar la constatación y devolver el resultado al juez de Capital.
    \end{example}

    \item \textbf{Notificaciones en domicilio real:} realizadas por el oficial notificador en el domicilio de las personas. Son actos delegados por el tribunal, típicamente fuera de su sede.

    \item \textbf{Actuaciones del oficial de justicia:} constataciones de inmuebles, órdenes de lanzamiento, embargo de bienes muebles en domicilios, etc.
\end{enumerate}

\begin{definition}{Ley 22.172}
Convenio interjurisdiccional de asistencia recíproca entre jueces de distintas jurisdicciones. Regula el modo de comunicarse decisiones judiciales entre jueces de distintas provincias o entre el fuero federal y el provincial.
\end{definition}

\section{Tiempo de los Actos Procesales --- Horarios Hábiles}

\begin{itemize}
    \item Para actos realizados en la \textbf{sede del tribunal}: el horario hábil es de \textbf{7:30 a 13:30 horas} (establecido por acordada de la Corte).

    Si una audiencia está en curso y se prolonga más allá de las 13:30, se habilita el horario sin necesidad de resolución expresa, mientras la audiencia no demande muchas más horas. De lo contrario, el juez puede suspenderla y continuarla al día siguiente.

    \item Para actos realizados \textbf{fuera de la sede del juzgado} (notificaciones, actuaciones del oficial de justicia): el horario hábil se amplía hasta las \textbf{21:00 horas}.

    Si el juez dispone la habilitación horaria (por alguna circunstancia especial), el oficial puede cumplir la diligencia fuera del horario hábil.
\end{itemize}

% ============================================================
% CAPÍTULO 3: PLAZOS PROCESALES
% ============================================================
\chapter{Plazos Procesales}

Los plazos son el tiempo dentro del cual es preciso cumplir cada acto procesal. El Código establece distintos plazos, que son los tiempos fijados para el cumplimiento de los actos procesales.

\section{Clasificación de los Plazos Procesales}

\subsection{Por su Origen}

\begin{table}[H]
\centering
\begin{tabularx}{\textwidth}{>{\bfseries\raggedright\arraybackslash}p{0.28\textwidth} >{\raggedright\arraybackslash}X}
\toprule
\textbf{Tipo} & \textbf{Descripción} \\
\midrule
Plazos Legales & Establecidos directamente por la ley (procesal o de fondo). Ejemplo: 5 días para apelar; 15 días hábiles para contestar la demanda. \\[0.4em]
Plazos Judiciales & Fijados por el juez ante la ausencia de plazo legal. Ejemplo: el juez fija al perito un plazo de 15 o 20 días para presentar su dictamen. \\[0.4em]
Plazos Convencionales & Los que las partes acuerdan entre sí. Las partes pueden pactar plazos diferentes a los legales si están de acuerdo. \\
\bottomrule
\end{tabularx}
\caption{Clasificación de los plazos procesales por su origen}
\end{table}

\subsection{Plazos Perentorios / Fatales / Preclusivos}

Nuestro ordenamiento procesal adopta el \textbf{sistema de la preclusión}: los plazos procesales, una vez vencidos, hacen caducar la potestad de ejercer el acto de que se trata.

\begin{note}
Los términos \textbf{perentorio}, \textbf{preclusivo} y \textbf{fatal} son sinónimos en la práctica forense. Son plazos que, una vez cumplidos, hacen perder el derecho no ejercido, sin necesidad de petición de la contraparte.
\end{note}

\begin{example}{Preclusión por vencimiento de plazo}
La parte tiene 5 días para plantear la apelación desde que es notificada. Si pasan 6 días y no apeló, perdió el derecho: precluyó la posibilidad. El mismo principio aplica a los 15 días hábiles para contestar la demanda: pasado ese plazo, la parte pierde el derecho sin que nadie tenga que pedirlo.
\end{example}

\subsection{Plazos Prorrogables e Improrrogables}

Hay plazos que, antes de su vencimiento, pueden ser objeto de \textbf{prórroga a petición de parte}. Ejemplo: un plazo judicial de intimación de 10 días (no fijado por ley) es revisable por el juez.

Otros plazos son absolutamente \textbf{improrrogables (fatales)}: el plazo de 5 días para apelar, por ejemplo.

\subsection{Plazos Ampliados por Distancia (Art. 158 CPCCN)}

\begin{formula}{Art. 158 CPCCN --- Ampliación por Distancia}
Cuando la persona notificada vive lejos de la sede del tribunal, los plazos legales se amplían: \textbf{un día más por cada 200 km} entre el lugar de notificación y la sede del tribunal. Si la fracción restante es \textbf{superior a 100 km}, se agrega un día adicional; si es \textbf{inferior a 100 km}, no se agrega.
\end{formula}

\begin{example}{Cálculo de ampliación por distancia}
\begin{itemize}
    \item Persona notificada a \textbf{470 km} del tribunal: 400 km = 2 días adicionales. La fracción restante (70 km) es inferior a 100 km, por lo que no se agrega otro día. \textbf{Total: 2 días de prórroga.}
    \item Persona notificada a \textbf{520 km}: 400 km = 2 días + fracción de 120 km (mayor a 100 km) = 1 día más. \textbf{Total: 3 días de prórroga.}
\end{itemize}
\end{example}

\subsection{Plazos Individuales y Plazos Comunes}

La generalidad de los plazos procesales son \textbf{individuales}: corren para cada parte desde su propia notificación.

Los \textbf{plazos comunes} corren para todas las partes desde la última notificación. Ejemplos: el plazo para alegar y (antes de la reforma) el plazo de prueba.

% ============================================================
% CAPÍTULO 4: TEORÍA DE LA COMUNICACIÓN Y NOTIFICACIONES
% ============================================================
\chapter{Teoría de la Comunicación y Notificaciones}

\section{Traslados y Vistas}

Son providencias por las cuales el juez decide poner en conocimiento de una parte la petición formulada por otra.

\begin{definition}{Traslado}
Se corre traslado cuando una parte hace una petición ante el juez y éste pone en conocimiento de las otras partes esa petición, para que ejerzan su derecho de defensa. El \textbf{traslado de la demanda} es el traslado por antonomasia.
\end{definition}

\begin{definition}{Vista}
Se llama `vista' a los traslados que se dan a los funcionarios del \textbf{Ministerio Público} (fiscal, defensor de menores e incapaces, representante del fisco). Hoy, entre partes todo son traslados, y cuando interviene un ministerio público se usa la denominación `vista'.
\end{definition}

\section{Oficios}

Son comunicaciones que se libran a otros jueces, o a funcionarios (secretarios, ministros del Poder Ejecutivo, incluso al Presidente de la Nación), y también a reparticiones públicas o privadas.

\begin{table}[H]
\centering
\begin{tabularx}{\textwidth}{>{\raggedright\arraybackslash}p{0.32\textwidth} >{\raggedright\arraybackslash}p{0.30\textwidth} >{\raggedright\arraybackslash}X}
\toprule
\textbf{Destinatario} & \textbf{¿Quién firma?} & \textbf{Norma} \\
\midrule
Presidente, ministros, secretarios de Estado; otros jueces de la misma circunscripción & El juez & Art. 38 CPCCN \\[0.4em]
Reparticiones públicas o privadas (bancos, empresas, organismos, etc.) & El secretario del tribunal (con transcripción de la providencia del juez) & Art. 38 CPCCN \\[0.4em]
Jueces de otras jurisdicciones (otras provincias) & Secretario/Juez con formato especial & Ley 22.172 \\[0.4em]
Jueces de otros países & Vía diplomática (Ministerio de Relaciones Exteriores) & Exhorto diplomático \\
\bottomrule
\end{tabularx}
\caption{Tipos de oficios según destinatario y firmante}
\end{table}

\begin{example}{Exhorto diplomático}
Para obtener los movimientos de una cuenta bancaria suiza de alguien investigado en causa penal, o para hacer declarar un testigo que vive en Francia o Perú: se libra un exhorto diplomático con colaboración del Ministerio de Relaciones Exteriores, que lo hace llegar al juez extranjero, quien cita al testigo y devuelve el resultado al juez argentino de origen.
\end{example}

\section{Principio General: Notificación por Ministerio de la Ley}

\begin{formula}{Art. 133 CPCCN --- Notificación por Ministerio de la Ley}
Todas las providencias se notifican a las partes por ministerio de la ley los \textbf{días martes y viernes} de la semana. El solo transcurso del martes o viernes siguiente al dictado de la resolución produce la notificación para todas las partes y sus auxiliares.
\end{formula}

La notificación ministerial se basa en la presunción de que los días martes o viernes los expedientes están a disposición de las partes para ser compulsados. Si un martes o viernes es feriado, la notificación se cumple el martes o viernes siguiente a ese feriado.

\subsection{La Nota de Asistencia}

Si el expediente no estaba disponible para ser consultado (p. ej., estaba `a despacho'), la parte puede dejar constancia de ello (antes: en el libro de asistencia físico; hoy: virtualmente a través del sistema). En ese caso, la notificación por ministerio de la ley \textbf{no se cumple ese día}.

\section{Notificación Personal o por Cédula (Art. 135 CPCCN)}

La notificación ministerial reconoce importantes excepciones: los actos procesales más trascendentales deben notificarse de modo \textbf{personal o por cédula}.

\textbf{Modalidades posibles:}

\begin{itemize}
    \item \textbf{Personalmente:} la parte concurre al expediente y deja una nota manifestando que se notifica de la resolución.
    \item \textbf{Por cédula:} en formato papel (cuando la parte no tiene domicilio electrónico) o electrónica (cuando ya constituyó domicilio electrónico).
    \item \textbf{Por acta notarial:} realizada por escribano.
    \item \textbf{Por telegrama colacionado o carta documento:} para la mayoría de los actos (menos traslado de demanda y notificación de sentencia).
\end{itemize}

\subsection{Actos que Requieren Notificación Personal o por Cédula (Art. 135 CPCCN)}

\begin{formula}{Art. 135 CPCCN --- Actos que requieren cédula}
Se notificarán por cédula o personalmente:
\begin{itemize}
    \item La resolución que dispone el traslado de la demanda y la reconvención, junto con los documentos que se acompañan.
    \item La que dispone el traslado de las excepciones previas.
    \item La que cita a absolver posiciones.
    \item La que declara la cuestión de puro derecho.
    \item La que cita a audiencia preliminar (art. 360).
    \item Las que ordenan intimaciones o apercibimientos no establecidos directamente por la ley.
    \item La que hace saber la devolución del expediente de cámara.
    \item El traslado de liquidaciones.
    \item La primera providencia luego de que el expediente vuelva del archivo o haya estado más de tres meses fuera de secretaría.
    \item La que cita a personas extrañas al proceso (testigos).
    \item Las sentencias definitivas y las interlocutorias con fuerza de tales.
    \item Las que ponen fin al proceso (p. ej., declaración de caducidad de instancia).
\end{itemize}
\end{formula}

\subsection{Casos en que la Cédula Sigue siendo en Formato Papel}

\begin{itemize}
    \item El traslado de la demanda (prácticamente siempre en papel, ya que el demandado aún no constituyó domicilio electrónico).
    \item La declaración de rebeldía del demandado que no compareció.
    \item La sentencia definitiva al demandado rebelde o que no compareció.
    \item Notificaciones a personas ajenas al proceso que no tienen domicilio electrónico constituido.
\end{itemize}

\subsection{El Procedimiento del Oficial Notificador en el Traslado de Demanda}

El oficial notificador va al domicilio real de la persona denunciado por la parte actora y sigue el siguiente procedimiento:

\begin{enumerate}
    \item Pregunta si la persona vive en ese domicilio.
    \item Si le informan que sí, deja un \textbf{aviso de que concurrió} y que volverá al día siguiente en un determinado horario.
    \item Al día siguiente vuelve: si la persona está, la notifica; si no está, \textbf{deja la cédula} porque le han informado que la persona vive allí.
    \item Si la persona \textbf{no vive} en ese domicilio, no puede dejar la cédula: informa lo que le manifestaron quienes viven allí y devuelve la cédula al juzgado, que espera instrucciones de la parte actora.
\end{enumerate}

\begin{note}
Existe además la posibilidad de \textbf{notificación bajo responsabilidad de la parte actora}: el oficial notifica aunque le digan que la persona no vive allí. Esta modalidad es frecuentemente atacable por vía de nulidad.
\end{note}

\section{Notificación por Edictos}

Se utiliza para notificar a:

\begin{itemize}
    \item \textbf{Personas inciertas:} por ejemplo, en los procesos sucesorios es obligatorio citar a posibles herederos o acreedores desconocidos del causante.
    \item \textbf{Personas de paradero desconocido:} cuando no es posible localizarlas en ninguno de sus domicilios registrados (RENAPER, padrón electoral, AFIP, etc.).
\end{itemize}

Los edictos se publican en \textbf{periódicos de amplia tirada} y en el \textbf{Boletín Oficial}. No siempre es un método efectivo para que la persona se entere, pero de alguna manera se cumple formalmente con la notificación.

% ============================================================
% CUADRO CORNELL
% ============================================================
\chapter{Cuadro de Repaso --- Método Cornell}

\begin{table}[H]
\centering
\begin{tabularx}{\textwidth}{
    >{\raggedright\arraybackslash}p{0.28\textwidth}
    >{\raggedright\arraybackslash}X
}
\toprule
\textbf{Preguntas / palabras clave} & \textbf{Notas / respuestas} \\
\midrule

\textbf{¿Qué es un acto procesal?} &
Hecho voluntario con efecto inmediato de constituir, desenvolver o extinguir un proceso judicial. Se distingue del hecho procesal, que opera por circunstancias ajenas a la voluntad (p.ej. muerte de una parte). \\[0.5em]

\textbf{¿Quiénes pueden realizar actos procesales?} &
Las partes y sus litis consortes; los auxiliares de las partes (letrados, procuradores); el juez y los auxiliares del tribunal (secretario, prosecretario, empleados); terceros vinculados (peritos, testigos, oficiales de justicia, martilleros, oficiales notificadores). \\[0.5em]

\textbf{¿Cuál es la clasificación de Palacio?} &
Actos de iniciación (demanda, contestación, diligencias preliminares); actos de desarrollo (instrucción: alegación y prueba; dirección: ordenación, transmisión, documentación, cautelares); actos de conclusión (sentencia y modos anormales). \\[0.5em]

\textbf{¿Qué son los actos de impulso?} &
Son los actos que hacen avanzar el proceso hacia la sentencia (siguiente eslabón en la cadena). Su importancia es crucial: son los únicos que interrumpen el plazo de caducidad de instancia. \\[0.5em]

\textbf{¿Cuáles son los modos anormales de extinción del proceso?} &
Desistimiento (actor, unilateral); allanamiento (demandado, unilateral); conciliación (bilateral, en el proceso); transacción (bilateral, fuera del proceso, contrato); caducidad de instancia (por inactividad, sanción procesal). \\[0.5em]

\textbf{¿Qué es la caducidad de instancia?} &
Extinción del proceso por falta de impulso durante el plazo legal (6 meses en juicio ordinario; 3 meses en ejecuciones e incidentes --- art. 310 CPCCN). Es una sanción y también un modo de finalización. Puede decretarse de oficio o a petición de la contraparte. \\[0.5em]

\textbf{¿Qué debe contener un escrito judicial?} &
SUMA (título); encabezamiento con nombre de la parte, nombre y tomo/folio del letrado, carátula del expediente, domicilio electrónico. Debe redactarse en tinta negra, en idioma nacional, conforme las acordadas de la Corte. \\[0.5em]

\textbf{¿Qué es el art. 120 CPCCN y cuál es su relevancia?} &
Obliga a presentar copia de todo escrito del que deba darse traslado. Hoy se cumple subiendo el archivo al sistema NextGen dentro de las 24 horas. El incumplimiento puede hacer tener el escrito por no presentado, con consecuencias gravísimas para el cliente. \\[0.5em]

\textbf{¿Qué es el cargo y para qué sirve?} &
Es la constancia de recepción del escrito en el tribunal: indica el tribunal, la fecha y si está firmado por letrado. Lo coloca generalmente el prosecretario administrativo. Le otorga fecha cierta al escrito y lo convierte en un acto público incorporado al expediente. \\[0.5em]

\textbf{¿Cuándo puede realizarse un acto procesal fuera del tribunal?} &
Cuando hay audiencias de parte o testigos que no pueden desplazarse; en reconocimientos judiciales de lugar; en actuaciones del oficial de justicia (notificaciones, embargos, lanzamientos); y mediante oficios (Ley 22.172 para otras jurisdicciones; exhortos diplomáticos para el exterior). \\[0.5em]

\textbf{¿Cuáles son los horarios hábiles?} &
En la sede del tribunal: de 7:30 a 13:30 horas. Fuera de la sede del juzgado (notificaciones, actuaciones del oficial): hasta las 21:00 horas. Puede habilitarse fuera de ese horario por disposición judicial. \\[0.5em]

\textbf{¿Cuáles son las diferencias entre plazos legales, judiciales y convencionales?} &
Legales: fijados por la ley (5 días para apelar; 15 días para contestar demanda). Judiciales: fijados por el juez ante ausencia de norma (plazo al perito). Convencionales: acordados entre las partes. \\[0.5em]

\textbf{¿Qué significa que un plazo sea perentorio/fatal?} &
Que al vencerse caduca el derecho de ejercer el acto, sin necesidad de petición de la contraparte. El sistema de preclusión impide ejercer el acto fuera de término. \\[0.5em]

\textbf{¿Cómo se amplían los plazos por distancia?} &
Según el art. 158 CPCCN: un día más por cada 200 km entre el lugar de notificación y la sede del tribunal. Si la fracción restante es mayor a 100 km, se agrega un día adicional. \\[0.5em]

\textbf{¿Cuál es la diferencia entre traslado y vista?} &
Los traslados son comunicaciones a las partes de peticiones de la contraria. Las vistas se dan a los funcionarios del Ministerio Público (fiscal, defensor de menores, representante del fisco). Ambas tienen la misma función: poner en conocimiento y solicitar pronunciamiento. \\[0.5em]

\textbf{¿En qué consiste la notificación por ministerio de la ley?} &
Todas las resoluciones se tienen por notificadas los días martes y viernes, por el solo transcurso de esos días hábiles, sin necesidad de notificación expresa. Se basa en la presunción de que el expediente está disponible para ser consultado esos días. \\[0.5em]

\textbf{¿Cuándo es obligatoria la cédula de notificación?} &
Para los actos más trascendentes del proceso, enumerados en el art. 135 CPCCN: traslado de demanda, reconvención, excepciones previas, citación a absolver posiciones, audiencia preliminar, intimaciones, devolución del expediente de cámara, citación de testigos, sentencias definitivas e interlocutorias con fuerza de tales, y resoluciones que pongan fin al proceso. \\[0.5em]

\textbf{¿Cuándo se usa la notificación por edictos?} &
Para notificar a personas inciertas (herederos o acreedores desconocidos en sucesiones) y a personas de paradero ignorado. Se publican en diarios de amplia tirada y en el Boletín Oficial. \\

\midrule

\multicolumn{2}{p{\textwidth}}{%
\textbf{Resumen:} La Unidad 7 de Derecho Procesal Civil aborda los actos procesales como hechos voluntarios que constituyen, desarrollan o extinguen el proceso judicial. Siguiendo la clasificación de Palacio, se agrupan en actos de iniciación (demanda, contestación, diligencias preliminares), actos de desarrollo (instrucción ---alegación y prueba--- y dirección ---ordenación, transmisión, documentación y cautelares---) y actos de conclusión (sentencia y modos anormales: desistimiento, allanamiento, conciliación, transacción y caducidad de instancia). Los actos de impulso son especialmente relevantes porque son los únicos que interrumpen la caducidad de instancia (art. 310 CPCCN: 6 meses en juicio ordinario, 3 meses en ejecuciones e incidentes). En materia de forma, los escritos deben contener suma, encabezamiento completo y domicilio electrónico; el cargo les otorga fecha cierta. Los plazos se clasifican por su origen (legales, judiciales, convencionales) y por su perentoriedad; el art. 158 CPCCN regula la ampliación por distancia. En la teoría de la comunicación, la notificación opera por ministerio de la ley (martes y viernes, art. 133 CPCCN), salvo para los actos del art. 135 que exigen cédula o notificación personal. Los edictos se reservan para personas inciertas o de paradero ignorado.
} \\

\bottomrule
\end{tabularx}
\end{table}

\end{document}
